\section{Introduction}\label{sec:introduction}

Let \(X\subset\mathbf P^4_{\C}\) be a smooth cubic threefold and let
\((J,\Theta)\) be its principally polarized intermediate Jacobian.  The
theta divisor has a unique singular point, an ordinary triple point at the
origin, and its projectivized tangent cone is \(X\)
\cite[Theorem~1]{BeauvilleTheta}.  Its blow-up
\[
 \sigma:M=\operatorname{Bl}_0\Theta\longrightarrow\Theta
\]
is smooth, with exceptional divisor \(X\) and normal bundle
\(\cO_X(-1)\).  Write \(i:\Theta\hookrightarrow J\), \(b=i\sigma\),
\(e:X\hookrightarrow M\), and
\[
 \Lambda=H^1(J,\Z),\qquad
 \theta=[\Theta]\in\bigwedge^2\Lambda.
\]
The principal polarization makes \(\Lambda\) a unimodular symplectic lattice
of rank ten.  We use the divided powers
\(\theta^{[r]}=\theta^r/r!\), which are integral.

The ordinary group \(H^3(\Theta,\Z)\) is \(\bigwedge^3\Lambda\) by singular
weak Lefschetz.  Blowing up the triple point adds ten generators, but their
placement relative to restriction to \(X\) and pushforward to \(J\) is not a
direct sum.  The following theorem describes that placement.

Define
\[
 L:\bigwedge^3\Lambda\longrightarrow\bigwedge^5\Lambda,
 \qquad L(\alpha)=\theta\wedge\alpha,
\]
and let \(S=\Sat(L\bigwedge^3\Lambda)\) inside
\(\bigwedge^5\Lambda\).  The cylinder correspondence of the universal line
on \(X\) and the principal polarization give an integral isomorphism
\[
 y:H^3(X,\Z)\xrightarrow{\sim}\Lambda;
\]
we fix its global sign in Section~\ref{sec:geometric-glue}.

\begin{theorem}[Integral middle lattice]\label{thm:main-lattice}
There is an exact sequence of free abelian groups
\begin{equation}\label{eq:main-exact}
 0\longrightarrow\bigwedge^3\Lambda
 \xrightarrow{\ b^*\ }H^3(M,\Z)
 \xrightarrow{\ e^*\ }H^3(X,\Z)\longrightarrow0.
\end{equation}
It splits as a sequence of abelian groups.  The Gysin image is
\begin{equation}\label{eq:main-saturation}
 b_*H^3(M,\Z)=S
 =L\bigwedge^3\Lambda+\theta^{[2]}\wedge\Lambda,
 \qquad S/L\bigwedge^3\Lambda\cong\Lambda/2\Lambda.
\end{equation}
More precisely, \((b_*,e^*)\) embeds \(H^3(M,\Z)\) in
\(S\oplus H^3(X,\Z)\), with image
\begin{equation}\label{eq:main-fibre-product}
 \left\{(s,\gamma):
 [s]=[\theta^{[2]}\wedge y(\gamma)]
 \text{ in }S/L\bigwedge^3\Lambda\right\}.
\end{equation}
In particular, \(H^3(M,\Z)\) is torsion-free of rank \(130\), and the
nonzero Smith factors of \(L\) are \(1^{110},2^{10}\).
\end{theorem}

The split sequence \eqref{eq:main-exact} alone does not see the arithmetic
content: any lift of \(\gamma\) has pushforward in the mod-two coset
\([\theta^{[2]}\wedge y(\gamma)]\).  Thus ``glue'' below always means the
simultaneous lattice \eqref{eq:main-fibre-product}, not a nonsplit extension
of abstract abelian groups.

The dual statement corrects another tempting reading of the ten factors of
two.  Put
\begin{equation}\label{eq:escape-definition}
 E_M=H^5(M,\Z)\big/
 \bigl(b^*H^5(J,\Z)+\operatorname{tors}\bigr).
\end{equation}

\begin{theorem}[Escape lattice]\label{thm:main-escape}
The group \(E_M\) is free of rank ten.  The image of
\(e_*H^3(X,\Z)\) in \(E_M\) is exactly \(2E_M\).  Consequently
\[
 \frac{H^5(M,\Z)}{b^*H^5(J,\Z)+e_*H^3(X,\Z)+\operatorname{tors}}
 \cong(\Z/2)^{10}.
\]
\end{theorem}

The lattice is intrinsic to the singular theta divisor in middle degree.

\begin{corollary}[Integral intersection cohomology]\label{cor:main-ih}
There are natural isomorphisms
\[
 IH^3(\Theta,\Z)\cong H^3(M,\Z)
\]
and an exact sequence
\[
 0\longrightarrow H^3(\Theta,\Z)
 \longrightarrow IH^3(\Theta,\Z)
 \longrightarrow H^3(X,\Z)\longrightarrow0.
\]
Under these identifications, Theorem~\ref{thm:main-lattice} computes the
integral lattice \(IH^3(\Theta,\Z)\).
\end{corollary}

The proof has three parts.  First, the link of the triple point is the circle
bundle of \(\cO_X(-1)\); its Gysin sequence and Mayer--Vietoris give
\eqref{eq:main-exact}.  Second, a symplectic-basis decomposition turns
\(L:\bigwedge^3\Lambda\to\bigwedge^5\Lambda\) into primitive blocks and
unsigned incidence matrices of complete graphs.  Each incidence block has
Smith form \(1^{n-1},2\), and its order-two class is exactly a divided-power
vector \(\theta^{[2]}\wedge z\).  Third, the difference map of the Fano
surface produces lifts of all ten quotient directions.  The identity
\[
 \xi\star\theta^{[3]}=\pm\theta^{[2]}\wedge y(\xi)
\]
for the Pontryagin product identifies their pushforwards and yields
\eqref{eq:main-fibre-product}.  No finite Smith-form computation enters the
proof.

Beauville established the singularity and degree-six Fano geometry used
here \cite{BeauvilleTheta}; Clemens and Griffiths supply the integral
cylinder isomorphism and the minimal class of the Fano surface
\cite{ClemensGriffiths}.  Krämer already computed the rational
decomposition-theorem correction for this resolution
\cite[Corollary~6]{KraemerE6}; Section~\ref{sec:intersection} explains why
that calculation does not determine the integral middle lattice.  The
integral Lefschetz filtration and compatible splittings of Faulkner Valiente
and Miller Eismeier already imply the Smith factors and the abstract
saturation quotient in \eqref{eq:main-saturation}
\cite[Theorem~2.9, Corollary~2.10, and Proposition~2.14]{FaulknerMiller}.
Our direct off-central calculation records the complete-graph blocks and the
specific representatives \(\theta^{[2]}\wedge z\) that the Fano endpoint
realizes geometrically in \eqref{eq:main-fibre-product}.

The stronger assertion that the full degree-six Fano transfer surjects onto
\(H^3(M,\Z)\) is not needed and is not claimed.  Only ten endpoint classes,
whose images in \(H^3(X,\Z)\) form a basis, enter the argument.
